\documentclass[11pt]{article}
\usepackage{amsmath,amssymb,amsthm}
\newtheorem{theorem}{Theorem}
\newtheorem{lemma}{Lemma}
\title{Butterfly Comparisons and Two-Word Paired Cauchy Readouts}
\author{}\date{}\begin{document}\maketitle

\paragraph{Status.}
These are complete written elementary arguments with exact finite checks,
not Lean proofs. They identify two available mechanisms. They do not
prove a growing-capacity rational decoder: the paired-root interpolation
below does not preserve any child-block constraints. No statement about
the unrestricted worst-case HeadComplexity follows from this note alone.

\begin{theorem}[All butterfly comparison patterns occur]
Consider $H=2^d$ wires with bits numbered $0,\ldots,d-1$, and $d$
layers of compare-exchange gates. Layer $k$ compares each pair differing
in bit $k$, sending the smaller value to the bit-zero output. Every
choice of the $Hd/2$ swap bits is the comparison pattern of some
assignment of $H$ distinct real input values. The values can be a
permutation of the fixed set $\{0,\ldots,H-1\}$.
\end{theorem}
\begin{proof}
Choose the swap pattern, and let $\pi$ be its routed input-to-output
permutation. For output label $y$, set
\[
 v(y)=\sum_{j=0}^{d-1} y_j2^{d-1-j}.
\]
Give input $h$ the value $v(\pi(h))$. At a layer-$k$ gate, its two
routed values have final output labels with the same lower $k$ bits
and opposite bit $k$: all later layers preserve bits $0,\ldots,k$.
In the order defined by $v$, the first differing bit read from low to
high decides the comparison. Thus the routed bit-zero output has the
smaller value at this gate, regardless of the two final higher-bit
strings. Inducting through the layers shows that the actual comparator
network follows exactly the prescribed routing. All input values are
distinct because $v$ and $\pi$ are permutations.
\end{proof}

\paragraph{Readout issue.}
This theorem supplies the full number of independent comparator bits.
It does not implement an intermediate comparison by a head. A rational
input-Haar sum generally includes values from many output-prefix groups;
extracting the selected group's comparison remains necessary.

\begin{lemma}[Cauchy interpolation]
Let $x_1,\ldots,x_m$ be distinct sample points and let
$p_1,\ldots,p_m$ be distinct poles disjoint from them. The matrix
$(1/(x_i-p_j))_{i,j}$ is invertible.
\end{lemma}
\begin{proof}
If a linear combination vanishes at all $m$ samples, clearing its
pole product gives a polynomial of degree at most $m-1$ with $m$
distinct roots, hence zero. Taking residues at the distinct poles
then makes every coefficient zero.
\end{proof}

\begin{theorem}[Two arbitrary root words with paired poles]
Fix $q$ distinct query points $t_1,\ldots,t_q$ and $2q$ distinct poles
$p_{h,0},p_{h,1}$, all disjoint from the queries. For any two sign words
$\sigma,\tau\in\{-1,1\}^q$, there are unrestricted real residues
$a_{h,0},a_{h,1}$ such that, writing
\[
 R_h(t)=\frac{a_{h,0}}{t-p_{h,0}}+
        \frac{a_{h,1}}{t-p_{h,1}},\qquad
 G_\rho(t)=\sum_{h=1}^q\rho_h R_h(t),
\]
one has $\sigma_iG_\sigma(t_i)>0$ and
$\tau_iG_\tau(t_i)>0$ at every query.
The statement holds in particular for positive exterior queries and
arbitrarily tightly paired positive poles below every query.
\end{theorem}
\begin{proof}
Replacing $\tau$ by $-\tau$ leaves its required inequalities unchanged,
because both the outside label and $G_\tau$ change sign. We may therefore
suppose that the agreement set $I=\{i:\sigma_i=\tau_i\}$ has size at
least $q/2$. Write $J=[q]\setminus I$.

The $2|I|$ poles belonging to groups in $I$ include at least $q$ distinct
poles. By Cauchy interpolation choose a rational sum $S_I$, using these
poles, with values $S_I(t_i)=\sigma_i$ for $i\in I$ and zero for
$i\in J$. If $J$ is nonempty, use $|J|$ of its $2|J|$ poles to
interpolate a sum $S_J$ with $S_J(t_j)=\sigma_j$ for $j\in J$.
If $J$ is empty take $S_J=0$.

Choose $L>\max_{i\in I}|S_J(t_i)|$. Assign the residues in the $I$
groups, after dividing by their fixed factors $\sigma_h$, so that
$\sum_{h\in I}\sigma_hR_h=L S_I$. Assign those in the $J$ groups so
that $\sum_{h\in J}\sigma_hR_h=S_J$. Then
\[
 G_\sigma=L S_I+S_J,\qquad G_\tau=L S_I-S_J.
\]
On $I$ the first term strictly dominates and has the desired common
sign. On $J$ it is zero, so the two values have the desired opposite
signs. This proves every strict inequality.
\end{proof}

\paragraph{Application and limitation.}
At the root of a butterfly, the two physical poles with lower prefix
$h$ contribute a difference of two simple reciprocals. Because the
residues in this theorem are unrestricted, its plus notation includes
that difference. Consequently any two prescribed root words can be
handled by one common residue vector, for fixed paired poles. Positivity
of the residues is not asserted. More importantly, the interpolation
uses the same residues that determine all lower input-Haar blocks and
need not respect their signs. The full-tree shared-cone problem is
therefore strictly stronger than this theorem.

\paragraph{The coalesced induction criterion.}
For completeness, consider fixed child strict feasibility cones
$K_0,K_1\subseteq\mathbb R^M$ in child-indexed residue coordinates.
When the newest pole pairs coincide, the two child residue vectors
$a,b$ contribute root values $C(a-b)$, where $C$ is the square Cauchy
matrix on the shared lower-prefix queries and poles. If the physical
root routing is $\theta_s$, the child coordinates are
$a_s=c_{s+\theta_sM}$ and $b_s=c_{s+(1-\theta_s)M}$; this accounts
for the swap only once. In particular the child cones are independent
of the root word at coalescence.

The root is feasible for every word if and only if
\[
 C(K_0-K_1)=\mathbb R^M,
 \quad\text{equivalently}\quad K_0\cap K_1\ne\varnothing. \tag{1}
\]
Indeed, a proper open convex cone is weakly separated from some point
by a nonzero linear functional, and then misses the open orthant
opposite to that functional. Thus meeting every orthant forces the
open convex cone to be all of $\mathbb R^M$. A cone containing zero
and open in $\mathbb R^M$ is the whole space; applying this to
$K_0-K_1$ proves the second equivalence. Empty child cones are excluded
from this discussion, since then no root can be feasible.

With a free root-block bias, replace the left side of (1) by
$C(K_0-K_1)+\operatorname{span}\{\mathbf1\}$. The corresponding
criterion is
\[
 (K_0-K_1)\cap\operatorname{span}\{C^{-1}\mathbf1\}\ne\varnothing.
 \tag{2}
\]
This is weaker than ordinary intersection, but still requires proof for
the actual hierarchy. Neither pairwise intersection for a smaller
hierarchy nor the two-root interpolation theorem supplies that proof:
two parent targets involve four child cones, with coordinatewise swaps.

\paragraph{Explicit bias direction.}
For $C_{ij}=1/(1-p_i-p_j)$, put
$P(t)=\prod_j(t-p_j)$ and $Q(t)=\prod_j(t-(1-p_j))$.
Then
\[
 (C^{-1}\mathbf1)_h=-\frac{Q(p_h)}{P'(p_h)},
\]
because the proper rational function $1-Q(t)/P(t)$ equals one at
every sample $1-p_i$. This vector has alternating signs when the
positive poles are ordered and all are less than $1/2$. Its masked
partial sums need not be constant on the child query blocks.

\paragraph{Exact verification.}
The companion checker exhausts all butterfly settings through eight
wires, checks the paired-root interpolation on all sign-word pairs
through four queries and on additional larger cases, and verifies the
explicit inverse-Cauchy constant vector. These finite checks do not
claim a growing-size full-tree decoder.
\end{document}
